\section{Introduction}\label{sec:introduction}

Hermitian--Yang--Mills flow on a degenerating curve can have several
relaxation scales. Long annuli weaken the interaction between normalization
components, while a component connection approaching a reducible connection
introduces further small endomorphism modes. An approximation on a fixed
multiple of the shortest annular length need not reach the time at which
the weakest mode relaxes. The problem is to compare actual trajectories
with a finite matrix evolution on the entire positive time axis, allowing
these two sources of small eigenvalues to vary independently.

We consider a compact genus-three plumbing family with two elliptic and
two rational normalization components. Four independent node parameters
control the annular lengths. Two independent real parameters control flat
unitary connections on the elliptic components. Fixed unitary identifications
at the nodes give a flat bundle on the entire family; its component and
cycle holonomies can be noncommuting. The finite configuration space is a
product of spaces of positive Hermitian matrices of determinant one. Its
potential is the sum of component commutator potentials and squared
distances across the nodes, with their geometrically determined weights.

The all-time comparison uses two different estimates for the error of the
reconstructed metric. The component and radial errors have an integrable
bound in terms of the finite forces. The first-order change from component
frames to parallel node frames has only the first angular Fourier modes.
A Neumann problem on each incidence collar writes this error as a
divergence whose space-time energy is small. The actual HYM energy identity
then controls its contribution to the trace comparison. This argument does
not divide by the smallest positive eigenvalue.

Throughout the all-time theorem, the component matrices and node
identifications are fixed. Only their nonnegative weights vary, and the
initial finite metrics lie in a specified small ball. Zero component
weights and nonzero common kernels are permitted. The approximation is an
absolute integrated squared-distance estimate; it does not compare errors
relative to a trajectory that is tending to zero.

\subsection{The main theorem}
Let $C_\tau$ be the smooth fibre of the four-node family constructed
in Section~\ref{sec:geometry}. Its normalization components have
area one. On the elliptic components the reference connections are
$d-i\alpha_iT_i\vartheta_i$, where $T_i$ are fixed trace-free
Hermitian matrices and $\vartheta_i$ fixed nonzero real harmonic
one-forms. The rational component connections are trivial.
The parameters $\alpha_i$ are real and independent of the four
complex node parameters $\tau_e\ne0$. Write $h_0$ for the
global parallel bundle metric and set
\[
 L_e=\log(1/|\tau_e|),\quad L_*=\min_eL_e,\quad
 \epsilon_e=L_e^{-1},\quad
 \alpha_*=\max_i|\alpha_i|,\quad
 \kappa_i=\int_{E_i}|\vartheta_i|^2.
\]
The joint limit is $\alpha_*\to0$ and $L_*\to\infty$.
All parameters are held constant during each flow.

Let $\cP_r$ be the positive determinant-one Hermitian matrices,
with metric $\langle V,W\rangle_H=\tr(H^{-1}VH^{-1}W)$
and distance $d$. Equip $\cP_r^4$ with the product metric
$g_A$ weighted by the component areas $A_i=1$.
The finite evolution is the negative gradient flow of
\begin{equation}\label{eq:intro-potential}
 \Phi(H)=
 \sum_{i=+,-}\alpha_i^2\kappa_i
   \{\tr(T_iH_i^{-1}T_iH_i)-\tr(T_i^2)\}
 +\pi\sum_{e=ij}\epsilon_e d(H_i,U_{ij}H_jU_{ij}^*)^2.
\end{equation}
Its matrix equation is printed in \eqref{eq:finite-flow}.
The determinant-one actual HYM equation is
$h^{-1}\partial_sh=-2i\Lambda F_h$ in physical time $s$.
For two metrics define
$D_{\tr}(h,k)=\tr(h^{-1}k+k^{-1}h-2I)$.

\begin{theorem}\label{thm:intro}
Fix the rank, component metrics, matrices $T_i$, harmonic forms
$\vartheta_i$ and unitary node identifications $U_e$. There
is $R_*>0$ and an explicit smooth reconstruction
$H\mapsto\widehat K_{\tau,\alpha}(H)$ such that the following
holds for sufficiently small joint parameters.
Let $H(s)$ be the finite gradient flow from distance at most
$R_*$ from $(I,I,I,I)$. Let $h(s)$ be actual HYM flow on
the same smooth fibre, from any smooth metric with the parallel
determinant and fixed bounds $ch_0\le h(0)\le C_0h_0$.
Then
\begin{equation}\label{eq:intro-main}
 \sup_{s\ge0}\int_{C_\tau}
     d(h(s),\widehat K_{\tau,\alpha}(H(s)))^2\,dA_\tau
 \le\eta+C\left(\alpha_*+
                   \frac{(1+\log L_*)^2}{L_*}\right),
\end{equation}
where
$\eta=\int_{C_\tau}D_{\tr}(h(0),\widehat K_{\tau,\alpha}(H(0)))\,dA_\tau$.
Constants are independent of all $\alpha_i^2L_e$,
all $L_e/L_f$, and time. Component weights may vanish and
the common space of parallel trace-free endomorphisms may
be nonzero. On each fixed compact normalization region the
reconstruction differs from $H_i(s)$ by $O(L_*^{-1})$.
\end{theorem}

The proof is Theorem~\ref{thm:all-time}. The exponentially
small geometric-transfer term in \eqref{eq:all-time-error}
is absorbed in the error displayed here. Uniform approximation
requires $\eta\to0$ over the chosen initial family.
Theorem~\ref{thm:spatial-entry} removes the need to prescribe
the finite initial tuple. It gives a nonempty open class of
spatially varying smooth metrics, defined by the component
extension norm \eqref{eq:extension-norm}, which enters the
comparison at a common time $s_*\to\infty$ satisfying
$(\alpha_*^2+L_*^{-1})s_*\to0$. The selected tuple depends
on the initial component profiles. Smooth middle-neck data
need only obey the fixed pointwise bound; their derivatives
need not be bounded uniformly.

Theorem~\ref{thm:relative-spectrum} identifies every small
eigenvalue of the actual endomorphism connection Laplacian.
Its finite energy is the sum of
$\alpha_i^2\kappa_i\|[T_i,X_i]\|^2$ and
$(2\pi/L_e)\|X_i-\Ad(U_e)X_j\|^2$, with mass
$\sum_i A_i\|X_i\|^2$. The comparison is relative for
every positive weak eigenvalue, identifies the exact kernel,
and leaves a uniform complementary gap. In contrast to the
all-time nonlinear constants, this spectral result is uniform
also when the transports vary through reducible tuples.

The same family contains a stable rank-two example with
three node weights $h$, a fourth node weight $h^2$, and
component weight $\alpha_+^2\kappa_+=h^4$. Its twelve
small eigenvalues comprise nine order-$h$ modes, two
asymptotic to $\tfrac12h^2$, and one asymptotic to $h^4$.
Corollary~\ref{cor:three-scales} proves both these spectral
claims and actual nonlinear approximation through $h^{-4}$
and all later times. The latter conclusion uses the all-time
residual argument, independently of the spectral asymptotics.

\subsection{The analytic mechanism}
The reconstruction uses scalar component Green functions,
exactly flat geodesic node profiles and relative-logarithm
interpolation on collars of radius $\rho=L_*^{-1}$.
The frame change on an elliptic incidence is
$e^{i\alpha_iT_iF_{ie}}$ with $F_{ie}$ a real linear
harmonic function in the branch coordinate. Estimating its
cutoff error in absolute $L^1$ can produce an integral
of size $1/|\alpha_i|$; the explicit calculation
\eqref{eq:first-order-divergence} explains this difficulty.

Instead, the term linear in $\alpha_i$ in the full Chern
expression has only angular Fourier indices $1,-1$ after
the conformal area factor is removed. Its entrywise zero
mean allows a collar Neumann solve with zero normal trace
at both boundaries. Extension by zero gives a global
divergence without a boundary measure. Thus the residual is
$D^*J+B$, with $J$ small in space-time $L^2$ and $B$ small
in space-time $L^1$. The trace comparison pairs $J$ with
derivatives of the actual and reconstructed metrics; both
have uniformly bounded total energy.

For the remainder $B$, a fixed-matrix weighted normal
equation estimate controls the separate finite forces by
the gradient on a small ball. Ohta's bounded-image
rectifiability theorem for self-contracted curves
\cite{ohta} then gives their uniform total integral.
The fixed-constraint restriction enters at this step.
The matrix counterexample in Section~\ref{sec:finite}
shows why a uniform lower bound for the full unweighted
matrix alone does not suffice.

\subsection{Relation to earlier work}
Burger's theorem gives relative scalar graph asymptotics
on hyperbolic surfaces with unrestricted ratios among
the short geodesic lengths \cite{burger}. Ji--Wentworth
prove spectral convergence for conic degenerations
\cite{ji-wentworth,ji-wentworth-correction}; their
correction replaces a false isoperimetric lemma, which
is not used here. Daskalopoulos--Wentworth establish
connection-valued spectral convergence and uniform
Yang--Mills estimates away from accidental reducibility
\cite[Theorem~4.13, Corollary~4.14 and
Proposition~5.8]{dw}. Our normalization components carry
reducible holonomy, so their uniform-gap argument does
not supply the additional slow dynamics considered here.

Haiden--Katzarkov--Kontsevich--Pandit develop
finite-dimensional asymptotic descriptions of HYM flow
on a fixed base \cite{hkkp}. Their Proposition~4.5
also compares exact and approximate metrics under an
integrable pointwise bound on the forcing. The
divergence estimate here uses spatially integrated
energy to treat a first-order collar term not controlled
uniformly under degeneration by that pointwise criterion.
Lam's fixed-base asymptotic
theorem for semistable bundles assumes a locally free
graded bundle, with additional rationality hypotheses
in its detailed reduction \cite{lam}. Applying these
results fibre by fibre does not give constants uniform
under pinching.

Previous work on nodal curves gives actual
noncommuting HYM approximation on bounded slow-time
intervals \cite{wuebben-effective} and treats
nontrivial gluing, spatial entry and all-time
comparison under a positive normalized rate
\cite{wuebben-gluing}. The present result removes
that rate requirement for fixed constraints on a
small basin and couples independent node and
component-holonomy scales in one compact family.
The new all-time step is the global divergence
decomposition of the full Chern residual; the
finite graph operator or the scalar multiscale
principle alone is not that step.
